Theorem 2 (Connection to the Shapley value). Letφ(i) denote the Shapley value [56] of an input variablei. Then, the Shap-
ley valueφ(i) can be represented as a weighted sum of causal effects involving the variablei, i.e.,φ(i) =∑
S⊆N\{i}
1
|S|+1wS∪{i}.
In other words, the effect of a causal pattern withm variables should be equally assigned to them variables in the computation
of Shapley values.
• Proof: By the definition of the Shapley value, we have
φ(i) =E
m
E
S⊆N\{ i}
|S|=m
[v(xS∪{i})−v(xS)]
= 1
|N|
|N|−1∑
m=0
1(|N|−1
m
)
∑
S⊆N\{ i}
|S|=m
[
v(xS∪{i})−v(xS)
]
= 1
|N|
|N|−1∑
m=0
1(|N|−1
m
)
∑
S⊆N\{ i}
|S|=m
∆v{i}(xS)
= 1
|N|
|N|−1∑
m=0
1(|N|−1
m
)
∑
S⊆N\{ i}
|S|=m

∑
L⊆S
wL∪{i}

 // by Theorem 5
= 1
|N|
∑
L⊆N\{i}
|N|−1∑
m=0
1(|N|−1
m
)
∑
S⊆N\{ i}
|S|=m
S⊇L
wL∪{i}
= 1
|N|
∑
L⊆N\{i}
|N|−1∑
m=|L|
1(|N|−1
m
)
∑
S⊆N\{ i}
|S|=m
S⊇L
wL∪{i} // sinceS⊇L ,|S| =m≥|L| .
= 1
|N|
∑
L⊆N\{i}
|N|−1∑
m=|L|
1(|N|−1
m
)·
(
|N|−|L|− 1
m−|L|
)
·wL∪{i}
= 1
|N|
∑
L⊆N\{i}
wL∪{i}
|N|−|L|−1∑
k=0
1(|N|−1
|L|+k
)·
(
|N|−|L|− 1
k
)
| {z }
αL
Then, we leverage the following properties of combinatorial numbers and the Beta function to simplify the term wL =∑|N|−|L|−1
k=0
1
(|N|−1
|L|+k)·
(|N|−|L|−1
k
)
.
(i) A property of combinitorial numbers.m·
(n
m
)
=n·
(n−1
m−1
)
.
(ii) The definition of the Beta function. Forp,q >0, the Beta function is defined as B(p,q ) =
∫ 1
0 xp−1(1−x)1−qdx.
(iii) Connections between combinitorial numbers and the Beta function.
◦ Whenp,q∈ Z+, we haveB(p,q ) = 1
q·(p+q−1
p−1 ).
◦ Form,n∈ Z+ andn>m , we have
(n
m
)
= 1
m·B(n−m+1,m).
19
